\chapter{Creer en ti más que tú mismo}

\epigraph{La fidelidad no nace de una prohibición. Nace de una fe.}{}

Vamos con la segunda palabra difícil. Si en el capítulo anterior hablamos de la indisolubilidad y descubrimos que, mirada desde el lado correcto, es el regalo más grande que alguien puede hacerte, ahora toca hablar de algo que suele provocar todavía más resistencia: la fidelidad.

Y es que la fidelidad tiene mala prensa. No entre quienes ya llevan años casados y han descubierto lo que significa. Sino entre quienes la miran desde fuera, desde antes, desde la cultura que nos rodea. Para muchos, fidelidad es sinónimo de restricción. De renuncia. De jaula.

El argumento suele formularse más o menos así: «¿Me estás diciendo que no puedo estar con nadie más? ¿En toda mi vida? ¿Nunca?» Dicho así, suena a condena. Como si el matrimonio fuera una puerta que se cierra ---para siempre--- sobre un mundo de posibilidades. Como si casarse fuera aceptar una limitación que nadie en su sano juicio querría.

Y si la fidelidad fuera eso ---una prohibición sexual disfrazada de virtud---, tendrían razón en resistirse. Nadie quiere vivir encerrado. Nadie quiere que le digan lo que no puede hacer durante el resto de su vida.

Pero ¿y si la fidelidad no fuera eso? ¿Y si fuera exactamente lo contrario?

\section{Una palabra que viene de la fe}

Hagamos algo que siempre ayuda: ir a la raíz. La palabra \textit{fidelidad} viene del latín \textit{fidelitas}, que a su vez viene de \textit{fides}. Y \textit{fides} no significa prohibición, ni restricción, ni renuncia. \textit{Fides} significa \textbf{fe}.\footnote{El término latino \textit{fides} abarca un campo semántico amplio: fe, confianza, lealtad, crédito, palabra dada. En el derecho romano, la \textit{fides} era el fundamento de todo pacto: la confianza en que el otro cumplirá lo acordado. Cf. Benveniste, É., \textit{Le vocabulaire des institutions indo-européennes}, vol. 1, Minuit, 1969.}

Fidelidad es, literalmente, un acto de fe. No fe en Dios ---aunque también---, sino fe en una persona. En ti.

Cuando alguien te es fiel, lo que te está diciendo no es «me resigno a no estar con otros». Lo que te está diciendo es: «Creo en ti. Creo en lo que eres. Creo en lo que puedes llegar a ser. Y apuesto mi vida entera a esa creencia.»

Es una declaración de fe, no un acto de renuncia.

¿Y si fuera al revés?

¿Y si, en lugar de mirar la fidelidad como algo que tú \emph{das} ---renunciando, sacrificándote, privándote---, la miráramos como algo que tú \emph{recibes}? Porque cuando tu mujer te es fiel, o cuando tu marido te es fiel, lo que está ocurriendo no es que se estén privando de algo. Lo que está ocurriendo es que alguien está depositando en ti una fe absoluta. Una fe que dice: «De entre todos los seres humanos del mundo, yo creo que tú eres la persona con la que quiero recorrer la vida. Y creo que tú eres capaz de ser el hombre, la mujer, que yo he visto en ti.»

Eso no es una jaula. Eso es el mayor voto de confianza que un ser humano puede recibir.

\section{La fe que te hace mejor}

Piensa un momento en las personas que han creído en ti a lo largo de tu vida.

Quizá fue un profesor que vio algo en ti que tú no veías. Que te dijo, cuando estabas a punto de tirar la toalla: «Tú puedes con esto.» Y algo cambió. No porque de repente fueras más listo o más capaz. Sino porque alguien te miró con ojos de fe, y esa mirada te hizo crecer.

Quizá fue un entrenador que te puso en el equipo cuando nadie más apostaba por ti. O un padre, o una madre, que nunca dejó de confiar en ti, ni siquiera cuando tú mismo habías dejado de hacerlo. O un amigo que, en tu peor momento, te dijo: «Yo sé quién eres de verdad. Y esto no eres tú.»

Todos hemos tenido, si hemos tenido suerte, al menos a una persona así. Alguien que creyó en nosotros más de lo que nosotros creíamos en nosotros mismos. Y todos sabemos lo que eso produce: no quieres fallarle. No porque te castigue si lo haces. No por miedo. Sino porque su fe en ti te eleva. Porque cuando alguien te mira así, traicionar esa confianza sería traicionarte a ti mismo.

Ahora multiplica eso por la relación más íntima de tu vida.

Cuando tu esposa te dice «sí, quiero», lo que está diciendo es: «Creo en ti con todo lo que tengo. Creo en ti con mi cuerpo, con mi futuro, con mis hijos, con mis años. Te entrego la fe más grande que puedo entregar: la fe de que tú vas a estar a la altura.» Cuando tu esposo te dice «sí, quiero», está haciendo exactamente lo mismo.

Eso no es una restricción. Es una responsabilidad enorme, sí. Pero es la clase de responsabilidad que te hace mejor persona. Porque la fe que alguien deposita en ti tiene un poder que nada más tiene: te muestra una versión de ti mismo que quizá no conocías. Y te hace querer ser esa versión.

\section{«No puedo» y «no quiero fallarte»}

Hay una diferencia abismal entre dos frases que, desde fuera, pueden parecer iguales.

La primera: «No puedo estar con otra persona.» La segunda: «No quiero fallarte.»

La primera es prohibición. Viene de fuera. Me la imponen. La obedezco ---si la obedezco--- a regañadientes, porque es la norma, porque es lo que toca, porque me metí en esto y ahora tengo que aguantar. La primera genera resentimiento. Acumula frustración. Convierte la fidelidad en una carga que pesa un poco más cada año.

La segunda es respuesta. Viene de dentro. Nace de haber recibido algo tan grande que no quiero destrozarlo. La segunda no necesita vigilancia ni control. No necesita que nadie me vigile, porque no es miedo lo que me mueve: es amor. Es la consciencia de que alguien me ha confiado lo más valioso que tiene ---su fe en mí--- y yo no quiero ser indigno de eso.

¿Ves la diferencia? En un caso, la fidelidad es una ley que me limita. En el otro, es una respuesta libre a un amor que me eleva. En un caso, soy un preso. En el otro, soy alguien que ha sido creído, y que elige estar a la altura de esa fe.

La fidelidad cristiana no funciona como una prohibición. Funciona como una correspondencia.\footnote{Benedicto XVI lo expresó con claridad al hablar del amor: «Él nos ha amado primero y sigue amándonos primero; por eso, nosotros podemos corresponder con el amor.» Benedicto XVI, Carta encíclica \textit{Deus Caritas Est}, n. 17 (2005). El mismo principio se aplica a la fidelidad: es respuesta a una fe recibida, no cumplimiento de una norma impuesta.} No digo «no» a otros porque me lo prohíban. Digo «sí» a ti, cada día, porque tu fe en mí merece mi mejor versión.

\section{La confianza como suelo firme}

Si habéis leído el capítulo 2, recordaréis lo que vimos sobre el apego. Bowlby descubrió que un niño necesita una \textit{base segura}: alguien predecible, constante, que esté ahí cuando lo necesite. Desde esa seguridad, el niño puede explorar el mundo. Sin ella, se paraliza o se desespera.\footnote{Cf. capítulo 2 de este libro, especialmente la sección sobre Bowlby y la teoría del apego.}

La fidelidad es la versión adulta, consciente y elegida de esa misma dinámica.

Cuando sabes que tu pareja te es fiel ---no porque le obligan, sino porque cree en ti y en lo vuestro---, ocurre algo parecido a lo que le ocurre al niño con apego seguro: puedes soltar la guardia. Puedes mostrarte como eres, con tus miserias, con tus contradicciones, con esas partes de ti que no enseñas a nadie. Puedes ser vulnerable sin miedo, porque sabes que el otro no va a usar tu vulnerabilidad como excusa para marcharse.

La biología ya nos inclinaba hacia ahí. Las hormonas del vínculo ---la oxitocina, la vasopresina--- nos empujan a la permanencia, a la confianza, al cuidado.\footnote{Cf. capítulo 3 de este libro, sobre la oxitocina y la vasopresina como hormonas del apego de pareja.} Pero la biología es un empujón, no una decisión. La fidelidad toma ese empujón biológico y lo convierte en promesa libre. Lo eleva de instinto a elección. De química a don.

Y ahí está la conexión con lo que vimos en el capítulo anterior: la indisolubilidad crea el espacio seguro; la fidelidad lo habita. La promesa de que no te vas a ir crea el marco; la fe que depositas en el otro cada día lo llena de contenido. No son dos cosas distintas. Son dos caras de lo mismo: un amor que dice «estoy aquí, creo en ti, y no me voy a ningún sitio».

\section{Lo que la infidelidad destruye de verdad}

Quizá ahora se entienda mejor por qué la infidelidad duele tanto. No es solo que duela el engaño, aunque duele. No es solo que duelan los celos, aunque duelen. Lo que de verdad destroza la infidelidad es que destruye la fe.

Cuando alguien te es infiel, lo que se rompe no es un acuerdo. Lo que se rompe es la creencia de que eras suficiente. La confianza de que el otro te miraba y veía a alguien por quien merecía la pena apostar la vida entera. Lo que se rompe es esa voz que decía «creo en ti» y que de repente se queda muda.

Por eso la infidelidad se siente como una traición ---porque lo es---. Pero no es solo una traición al pacto. Es una traición a la persona. A la fe que esa persona había puesto en ti. A la versión de ti mismo que el otro te ayudaba a ser.

Y por eso ---no por moralismo, no por rigidez--- la Iglesia la toma tan en serio.\footnote{«La fidelidad expresa la constancia en mantener la palabra dada. Dios es fiel. El sacramento del Matrimonio hace entrar al hombre y a la mujer en la fidelidad de Cristo a su Iglesia.» \textit{Catecismo de la Iglesia Católica}, n. 2365.} No porque le importe controlar la vida sexual de nadie, sino porque sabe lo que está en juego: la fe de una persona en otra persona. Y esa fe es sagrada.

\section{Creer en ti cuando tú no crees en ti}

Hay algo más, y quizá sea lo más hermoso de todo.

La fidelidad no consiste solo en creer en el otro cuando es fácil hacerlo. Cuando las cosas van bien, cuando estáis enamorados, cuando todo fluye. Eso no requiere mucha fe: es como creer que el sol saldrá mañana.

La verdadera fe se prueba cuando es difícil creer. Cuando el otro ha fallado. Cuando ha mostrado su peor versión. Cuando te ha decepcionado o se ha decepcionado a sí mismo. Cuando piensa que no vale la pena, que no es capaz, que no da la talla.

En esos momentos, la fidelidad es decir: «Yo sigo creyendo en ti. Aunque tú ahora mismo no creas en ti, yo sí. Yo sé quién eres. He visto lo mejor de ti, y sé que eso es lo real, aunque ahora esté escondido.»

Eso es fidelidad. No la fidelidad fácil de los días buenos, sino la fidelidad tozuda de los días malos. La que sostiene al otro cuando no puede sostenerse solo. La que le recuerda quién es cuando se le ha olvidado.

¿No es eso, en el fondo, lo que Dios hace con nosotros? ¿No es eso lo que significa que Dios sea fiel? Que sigue creyendo en nosotros cuando nosotros hemos dejado de creer. Que nos mira con ojos de fe incluso cuando nosotros nos miramos con ojos de vergüenza. Que su fidelidad no depende de nuestro merecimiento, sino de su amor.\footnote{«Si somos infieles, Él permanece fiel, porque no puede negarse a sí mismo» (2~Tim 2,13). La fidelidad de Dios no es una respuesta a nuestra bondad, sino una expresión de su ser. Y el matrimonio cristiano está llamado a reflejar, imperfectamente pero realmente, esa misma fidelidad. Cf. Juan Pablo II, Exhortación apostólica \textit{Familiaris Consortio}, n. 20 (1981).}

El matrimonio nos invita a algo extraordinario: a ser para el otro ese espejo de Dios. A creer en el otro como Dios cree en nosotros. No porque el otro sea perfecto ---no lo es, y tú tampoco---, sino porque la fe no se basa en la perfección del otro. Se basa en el amor que eliges tener.

\section{La fidelidad que se elige cada mañana}

Una última cosa. La fidelidad no es algo que se decide una vez, el día de la boda, y que luego funciona sola, como un piloto automático. La fidelidad se elige. Y se elige cada día.

Cada mañana, al despertar junto a la misma persona, hay una elección silenciosa que casi nunca se verbaliza pero que lo sostiene todo: «Hoy, otra vez, creo en ti. Hoy, otra vez, apuesto por nosotros.» No es una elección heroica. No sale en las noticias. Nadie te aplaude por ella. Pero es la elección más importante que vas a tomar en todo el día.

Y aquí está la paradoja hermosa de la fidelidad: cuanto más la eliges, más fácil se vuelve. No porque te acostumbres ---la costumbre mata---. Sino porque cada día que eliges creer en el otro, cada día que respondes a su fe con tu fe, el vínculo se hace más fuerte. La confianza se hace más honda. Y lo que empezó como una decisión se convierte, poco a poco, en una segunda naturaleza. En una forma de estar en el mundo. En la persona que eres.\footnote{Santo Tomás de Aquino describe la fidelidad como un hábito de la voluntad: no basta un acto aislado, sino la disposición estable a mantener la fe prometida. La virtud, en la tradición tomista, no es una imposición sino una perfección de la libertad. Cf. \textit{Summa Theologiae}, II-II, q. 110, a. 3.}

El día de la boda, la fidelidad es una promesa. Diez años después, es una historia. Treinta años después, es una identidad. No es que seas fiel porque prometiste serlo. Es que eres fiel porque en eso te has convertido, día a día, elección a elección, fe a fe.

Y eso ---convertirte en alguien fiel--- es mucho más grande que simplemente «no estar con otros». Es haberte transformado en la clase de persona que cree en otra persona con todo lo que tiene. Y eso, si lo piensas bien, se parece bastante a lo que significa amar.

\section{Porque debo, porque quiero, porque te quiero}

Hay una progresión que lo resume todo. Tres razones para hacer las cosas, y cada una supone un salto respecto a la anterior.

La primera: \textit{porque debo}. Es la razón de la norma, de la ley, del deber externo. Hago las cosas porque me las mandan, porque es lo que toca, porque las consecuencias de no hacerlas serían peores. Es el nivel más bajo. Cumples, pero no amas. Es como pagar impuestos: necesario, pero no te hace mejor persona.

La segunda: \textit{porque quiero}. Aquí hay un salto enorme. Ya no es imposición: es elección. Es la libertad en acto. Decía san Josemaría Escrivá que ``porque quiero'' es la razón más sobrenatural que existe,\footnote{Cf. capítulo 6 de este libro.} y tenía razón: en esas dos palabras se manifiesta la libertad que Dios nos ha dado, la misma que nos hace imagen suya. ``Porque quiero'' transforma el deber en don.

Pero hay una tercera, y es la que cambia las reglas del juego: \textit{porque te quiero}. Ya no es solo libertad ---es libertad orientada. Dirigida a una persona concreta. Con nombre, con cara, con defectos que conoces de memoria y virtudes que solo tú has visto. ``Porque te quiero'' es el amor que deja de ser abstracto y se encarna en alguien.

Y ahí está el misterio: ``porque te quiero'' se parece mucho a cómo ama Dios. Porque Dios no dice ``amo al mundo'' en genérico, como quien firma un manifiesto. Dios ama a cada uno de nosotros en particular, por nuestro nombre, de un modo irrepetible e insustituible.\footnote{«No temas, que yo te he rescatado, te he llamado por tu nombre y eres mío» (Is 43,1). El amor de Dios no es un amor de masas: es un amor personal, dirigido, con nombre propio.} El cónyuge hace exactamente lo mismo. No ama ``el amor''. No ama una idea. Te ama a ti. Con todo lo que eso implica.

La fidelidad recorre ese camino. Empieza en el ``debo'' de quien no conoce otra razón. Madura en el ``quiero'' de quien descubre la libertad. Y culmina en el ``te quiero'' de quien ha aprendido que el amor más alto no es el que se da a todos por igual, sino el que se entrega entero a una persona ---esta persona, aquí, ahora, para siempre---.

\paracaminar

\begin{enumerate}
  \item ¿Alguna vez alguien creyó en vosotros cuando vosotros no creíais en vosotros mismos? ¿Qué efecto tuvo esa fe? ¿Cómo os hizo sentir? Compartidlo.

  \item ¿Hay diferencia, en vuestra experiencia, entre el «no puedo» de la prohibición y el «no quiero fallarte» de la correspondencia? ¿Desde cuál de los dos vivís ---o queréis vivir--- la fidelidad?

  \item La fidelidad se elige cada día. ¿Qué gestos concretos y cotidianos os ayudan a renovar esa elección? ¿Qué podríais hacer para que vuestra fe mutua sea más visible, más tangible?
\end{enumerate}

\vinetacapitulo{ch09}
